\section{Relevância}
\label{sec:relevancia}

A chikungunya continua circulando no Brasil em ondas recorrentes \cite{desouza2023}, e no Recife a maior
parte da população ainda é suscetível \cite{braga2023}. Como a forma crônica da doença é frequente
\cite{paixao2018}, cada nova onda deixa um número grande de pessoas com dor articular por meses. Saber
quanta imunidade a cidade acumulou, em que ritmo a transmissão cresce e quando uma ação contra o mosquito
ainda faz diferença é útil para a vigilância e para o planejamento do controle vetorial.

Três argumentos justificam o projeto agora.

\textbf{Maturidade técnica.} Pela primeira vez há, para o Recife, uma série pública de notificações de
chikungunya por semana e por bairro, de 2015 a 2025 \cite{recife2026}, e um inquérito sorológico de base
populacional \cite{braga2023}. A série permite ajustar a dinâmica do modelo, e o inquérito fixa quantas
pessoas foram de fato infectadas, o que separa a transmissão da subnotificação. As ferramentas para
integrar e ajustar modelos de equações diferenciais são abertas e maduras \cite{rcore2025, soetaert2010}.

\textbf{Capacidade local.} Os resultados preliminares já incluem o tratamento reprodutível dos dados do
portal, um modelo SIR conferido contra todas as suas expressões analíticas e o modelo SEIR-SEI
implementado e verificado contra a teoria (Seção \ref{sec:preliminares}). O projeto começa com a base
pronta.

\textbf{Impacto esperado.} O modelo calibrado vai produzir estimativas da imunidade remanescente e cenários
de controle vetorial para a cidade, em código aberto que pode ser atualizado com as notificações de cada
ano. Sistemas como o InfoDengue já estimam a transmissão a partir das notificações \cite{codeco2018}. A
contribuição deste projeto é complementar: um modelo mecanístico, com o mosquito e a imunidade acumulada
representados de forma explícita, calibrado para o Recife e ancorado num inquérito sorológico local. Não
encontramos, na revisão feita para esta proposta, um modelo com essas características para a cidade.
